\documentclass[10pt,a4paper]{amsart}

\pdfinfoomitdate=1
\pdftrailerid{}
\pdfsuppressptexinfo=15

\usepackage[T1]{fontenc}
\usepackage{lmodern}
\usepackage[margin=1in]{geometry}
\usepackage{microtype}
\usepackage{amsmath,amssymb,mathtools}
\usepackage{booktabs}
\usepackage{natbib}
\usepackage[hidelinks]{hyperref}

\hypersetup{
  pdftitle={},
  pdfauthor={},
  pdfsubject={},
  pdfkeywords={}
}

\newtheorem{theorem}{Theorem}[section]
\newtheorem{proposition}[theorem]{Proposition}
\newtheorem{corollary}[theorem]{Corollary}
\newtheorem{lemma}[theorem]{Lemma}
\theoremstyle{remark}
\newtheorem{remark}[theorem]{Remark}
\newtheorem{definition}[theorem]{Definition}

\newcommand{\Exp}{\operatorname{Exp}}
\newcommand{\Nbh}{\mathcal N}

\title{Weighted Reverse-Stick Laws for Random Greedy Independent Sets on Threshold Graphs}
\author{Anonymous}
\date{}

\begin{document}

\begin{abstract}
The support of maximal independent sets of a threshold graph and the generic
random-greedy independent-set process are both known.  Assigning those facts
zero credit, we compute the endpoint law when vertices have arbitrary fixed
positive exponential rates.  In creation order, dominant vertices contribute
reverse hazards $h_d=w_d/W_d$, and the endpoint masses form a reverse
stick-breaking distribution.  The terminal law recovers every $h_d$ and gives
an explicit bijection between hazard vectors and the open simplex on the known
support, while leaving the original vertex rates nonidentifiable.  We derive
the accepted-update size PGF, all vertex marginals, and nested inclusion laws
for zero vertices.  Finally, we separate full-order inspections, accepted
updates, priority-label span, and continuous elapsed completion time; only the
last obeys a state-dependent Laplace recursion.
\end{abstract}

\maketitle

\section{Ownership, carrier, and process}

Random sequential adsorption on graphs is an established process
\citep{pippenger1989}, and random-order greedy maximal independent sets have a
large modern theory \citep{krivelevichetal2024}.  Threshold graphs and the
facets of their independent-set complexes are also known
\citep{klivans2007}.  Accordingly, this note claims no credit for the greedy
algorithm, its maximality, or the possible endpoint sets.  Its narrow residual
is the positive-rate endpoint distribution and the inverse and marginal
consequences of that distribution.

Let $v_1,\ldots,v_n$ be a threshold-graph creation order.  Set $b_1=0$; for
$i>1$, a creation bit $b_i=0$ makes $v_i$ isolated from all earlier vertices,
whereas $b_i=1$ makes $v_i$ adjacent to every earlier vertex.  Equivalently,
for $i<j$,
\begin{equation}\label{eq:adjacency}
 v_i\sim v_j\quad\Longleftrightarrow\quad b_j=1.
\end{equation}
Write
\[
 Z=\{i:b_i=0\},\qquad D=\{i:b_i=1\}.
\]

Give $v_i$ a fixed rate $w_i>0$ and an independent priority
$X_i\sim\Exp(w_i)$.  Scan vertices in increasing priority and accept a vertex
if no previously accepted neighbor blocks it.  Equivalently, among the active
vertices accept the first clock and delete its closed active neighborhood,
then continue.  The induced weighted order is the classical Plackett model
\citep{plackett1975}; that order and its exponential realization receive zero
credit here.  Let $I$ be the final independent set and
\begin{equation}\label{eq:prefix}
 W_j=\sum_{i=1}^j w_i.
\end{equation}

\section{Owned support}

The following is the threshold-graph maximal-independent-set support stated
by \citet{klivans2007}.  We include the one-line proof only to fix labels; it
is not a contribution of this note.

\begin{proposition}[Classical support; zero credit]\label{prop:support}
The maximal independent sets are exactly
\begin{equation}\label{eq:support}
 S_Z=Z,\qquad
 S_d=\{d\}\cup\{i>d:b_i=0\},\quad d\in D.
\end{equation}
\end{proposition}

\begin{proof}
Two dominant-creation vertices are adjacent.  If a maximal independent set
contains $d\in D$, it contains no earlier vertex and must contain every later
zero vertex; these vertices form $S_d$, while every omitted vertex is adjacent
to $d$ or to a later dominant vertex.  If it contains no dominant vertex,
maximality forces every zero vertex, giving $S_Z$; since $1\in Z$, every
dominant vertex has an earlier neighbor in $S_Z$.
\end{proof}

\section{Weighted endpoint law and inverse problem}

For $d\in D$ define the dominant prefix hazard
\begin{equation}\label{eq:hazard}
 h_d=\frac{w_d}{W_d}\in(0,1).
\end{equation}

\begin{theorem}[Weighted reverse-stick endpoint law]\label{thm:stick}
For every positive rate vector,
\begin{align}
 p_d:=\Pr(I=S_d)
 &=h_d\prod_{\substack{j\in D\\j>d}}(1-h_j)
 =\frac{w_d}{W_d}\prod_{\substack{j\in D\\j>d}}
       \frac{W_{j-1}}{W_j},                       \label{eq:pd}\\
 p_Z:=\Pr(I=S_Z)
 &=\prod_{j\in D}(1-h_j)
 =\prod_{j\in D}\frac{W_{j-1}}{W_j}.             \label{eq:pz}
\end{align}
These masses are positive and sum to one.
\end{theorem}

\begin{proof}
Read the creation string from the right.  If $b_n=0$, then $v_n$ is isolated
from the prefix and is always accepted; deleting it from the description
leaves the prefix endpoint law unchanged.  If $b_n=1$, then $v_n$ is universal.
It has the smallest priority among $v_1,\ldots,v_n$ with probability
$w_n/W_n=h_n$, in which case $I=\{n\}=S_n$.  Otherwise the first prefix vertex
is accepted and deletes $v_n$.  Conditional on this event, its identity has
probability $w_i/W_{n-1}$, and memorylessness leaves precisely the weighted
greedy law on the prefix.  Thus a terminal dominant bit takes an $h_n$ atom
and multiplies every prefix atom by $1-h_n$; a terminal zero bit only appends
that zero.  Iterating gives \eqref{eq:pd}--\eqref{eq:pz}.  The same recursion
proves positivity and normalization.
\end{proof}

The endpoint law does more than evaluate forward probabilities: it identifies
exactly the hazard coordinates that it can observe.

\begin{theorem}[Hazard inversion and open simplex]\label{thm:inverse}
From the labelled endpoint masses one recovers every dominant hazard by
\begin{equation}\label{eq:inverse}
 h_d=\frac{p_d}{p_Z+\displaystyle\sum_{\substack{e\in D\\e\leq d}}p_e},
 \qquad d\in D.
\end{equation}
Consequently, the map from $(h_d)_{d\in D}\in(0,1)^{|D|}$ to the masses in
\eqref{eq:pd}--\eqref{eq:pz} is a bijection onto the open probability simplex
on the $|D|+1$ sets in \eqref{eq:support}.  Every such mass vector is realized
by positive vertex rates.
\end{theorem}

\begin{proof}
Order $D$ increasingly.  Reverse-stick telescoping gives
\begin{equation}\label{eq:survival}
 p_Z+\sum_{\substack{e\in D\\e\leq d}}p_e
 =\prod_{\substack{j\in D\\j>d}}(1-h_j).
\end{equation}
Dividing \eqref{eq:pd} by \eqref{eq:survival} proves
\eqref{eq:inverse} and injectivity.

Conversely, start with strictly positive masses summing to one and define
$h_d$ by \eqref{eq:inverse}.  Its denominator strictly exceeds $p_d$, so
$0<h_d<1$.  Reversing the telescoping calculation recovers every prescribed
mass, proving surjectivity at the hazard level.  To realize the hazards by
rates, choose arbitrary positive rates at zero positions.  Moving from left
to right, when $d\in D$ set
\begin{equation}\label{eq:realize}
 w_d=\frac{h_d}{1-h_d}W_{d-1}.
\end{equation}
Since $b_1=0$, $W_{d-1}>0$; equation \eqref{eq:realize} is positive and makes
$w_d/W_d=h_d$.
\end{proof}

\begin{remark}[What is not identifiable]\label{rem:nonid}
The terminal law identifies the hazard vector, not the original rate vector.
All positive rate vectors satisfying $w_d/W_d=h_d$ at dominant positions are
observationally equivalent.  Global rescaling is always invisible.  More
strongly, the construction in \eqref{eq:realize} permits arbitrary positive
choices at every zero position and then adjusts later dominant rates.  Zero
rates after the last dominant position do not enter any hazard at all.  Thus
the simplex theorem is not a vertex-rate identifiability theorem.
\end{remark}

\section{Accepted-size PGF and inclusion laws}

Let $K$ be the number of accepted active-set updates.  Since every update adds
one vertex, $K=|I|$.  For $d\in D$ put
\begin{equation}\label{eq:kd}
 k_d=1+|\{i>d:b_i=0\}|.
\end{equation}

\begin{theorem}[Accepted-update PGF]\label{thm:pgf}
The complete PGF of $K$ is
\begin{equation}\label{eq:pgf}
 \mathbb E z^K=p_Zz^{|Z|}+\sum_{d\in D}p_dz^{k_d}.
\end{equation}
If distinct dominant positions give the same $k_d$, their masses are added at
that coefficient.  In particular,
\begin{align}
 \mathbb EK&=p_Z|Z|+\sum_{d\in D}p_dk_d,\label{eq:meanK}\\
 \operatorname{Var}(K)&=p_Z|Z|^2+\sum_{d\in D}p_dk_d^2-(\mathbb EK)^2.
 \label{eq:varK}
\end{align}
\end{theorem}

\begin{proof}
Proposition~\ref{prop:support} gives $|S_Z|=|Z|$ and $|S_d|=k_d$.
Mixing these deterministic sizes with Theorem~\ref{thm:stick} proves
\eqref{eq:pgf}; differentiation at $z=1$ proves the moment formulas.
\end{proof}

\begin{corollary}[Marginals and nesting]\label{cor:marginal}
For a dominant vertex $d\in D$ and a zero vertex $i\in Z$,
\begin{align}
 \Pr(d\in I)&=p_d,                                      \label{eq:dom-marginal}\\
 \Pr(i\in I)&=\prod_{\substack{j\in D\\j>i}}(1-h_j).
                                                               \label{eq:zero-marginal}
\end{align}
If $i,k\in Z$ and $i<k$, then
\begin{equation}\label{eq:nesting}
 \{i\in I\}\subseteq\{k\in I\},\qquad
 \Pr(i,k\in I)=\Pr(i\in I).
\end{equation}
\end{corollary}

\begin{proof}
Among the support sets, only $S_d$ contains dominant vertex $d$, proving
\eqref{eq:dom-marginal}.  A zero vertex $i$ is present precisely when no later
dominant vertex supplies the endpoint.  Summing $p_Z$ and the $p_d$ for $d<i$
or, equivalently, using \eqref{eq:survival}, gives
\eqref{eq:zero-marginal}.  Every support set containing an earlier zero $i$
also contains every later zero $k$, which proves \eqref{eq:nesting}.
\end{proof}

\section{Four count and clock objects}

The reverse-stick formula does not turn every implementation statistic into
the same random variable.  Under the fixed priorities $X_i$ define
\[
 J=n,\qquad K=|I|,\qquad
 R=\max_iX_i-\min_iX_i,\qquad
 \tau=\max_{v_i\in I}X_i.
\]
Here $J$ is the number of inspections in an implementation that scans the
entire priority order; $K$ is the number of accepted active-set updates and is
governed by \eqref{eq:pgf}; $R$ is the numerical span of all priority labels;
and $\tau$ is the continuous time at which the active set becomes empty.  The
last equality for $\tau$ follows in the one-shot priority coupling because the
last accepted clock deletes the last active vertices.  Neither $R$ nor $\tau$
is a discrete step count.

For an active vertex set $A$, let $\Nbh_A[v]$ be the closed neighborhood of
$v$ within $A$, and set $L_A(s)=\mathbb E_Ae^{-s\tau}$.  Exponential races give
the valid recursion
\begin{equation}\label{eq:ct}
 L_\varnothing(s)=1,\qquad
 L_A(s)=\frac{\displaystyle\sum_{v\in A}w_v
 L_{A\setminus\Nbh_A[v]}(s)}
 {s+\displaystyle\sum_{v\in A}w_v},\qquad s\geq0.
\end{equation}
Indeed, the first holding time is exponential with the denominator's total
rate, its winner has probability proportional to $w_v$, and memorylessness
restarts the residual active set.  Formula \eqref{eq:ct} is state dependent;
it is not obtained by substituting a Laplace variable into \eqref{eq:pgf}.

For a two-vertex clique with rates $1$ and $2$, the separation is already
literal: $J=2$, $K=1$, and at $s=1$
\begin{equation}\label{eq:firewall}
 \mathbb Ee^{-\tau}=\frac34,\qquad
 \mathbb Ee^{-R}=\frac13\frac23+\frac23\frac12=\frac59,
 \qquad \mathbb E(1/2)^K=\frac12.
\end{equation}
The numerical comparison is only a vocabulary firewall; the three transforms
belong to different statistics.

\section{Exact controls and scope}

The paper-local program independently enumerates the literal active-set chain
and uses rational arithmetic only.  It compares the resulting endpoint law
with Theorem~\ref{thm:stick}, then checks inversion, size PGFs, every marginal,
and every ordered pair of zero vertices.  A separate lane realizes prescribed
positive simplex masses, and another enumerates full weighted permutations to
check \eqref{eq:ct} through size six.

\begin{table}[htbp]
\centering
\caption{Exact-arithmetic falsification envelope.}\label{tab:controls}
\begin{tabular}{lr}
\toprule
control lane & inputs or checked cells\\
\midrule
parameter-labelled endpoint inputs & $31{,}833$\\
endpoint cells & $114{,}890$\\
accepted-size PGF cells & $65{,}404$\\
simplex realizations & $1{,}023$\\
continuous-time transform inputs & $567$\\
exact assertions & $750{,}181$\\
\bottomrule
\end{tabular}
\end{table}

The main endpoint grid contains every creation string through $n=6$ and all
weights in $\{1,2,3\}^n$, plus every creation string through $n=10$ under four
fixed uniform and nonuniform profiles.  Finite controls do not prove arbitrary
positive-real-rate claims and provide no novelty evidence.

Klivans's support, the generic RSA/MIS process, weighted orders, and
exponential races remain fully owned.  The defensible internal residual is
only Theorems~\ref{thm:stick}--\ref{thm:pgf} and
Corollary~\ref{cor:marginal}.  Because the right-to-left proof is short and may
be unpublished folklore, external circulation remains on hold pending a
deeper specialist owner search.

\bibliographystyle{plainnat}
\bibliography{references}

\end{document}
